\section{模解释与整数模型}
\label{sec:ch07-modular-interpretation}
% ============================================================================

上一节把 $X_0(N)$ 放回了 $\Q$ 上的代数几何世界，
这已经足够讨论有理点、Galois 作用与 Hecke 对应。
但如果我们的目标是第\ref{ch:eichler-shimura-l}章的 Eichler--Shimura 关系，
那么只停留在 $\Q$ 上仍然不够。
原因很简单：
Frobenius 元住在有限域里，
而有限域来自把整数模掉一个素数 $p$；
所以若想把复解析 Hecke 算子和算术 Frobenius 放到一起，
中间必须先有一座“从 $\Q$ 到 $\F_p$”的桥。
这座桥就是整模型。

\textbf{我们遇到了什么问题？}
一条定义在 $\Q$ 上的曲线，只告诉我们“有一组有理系数方程”。
可一旦把分母清掉并模掉一个素数，
得到的平面曲线很可能冒出奇点，
而群结构、同源乃至 Hecke 对应也可能在特殊纤维上失控。
因此真正要问的，不是“能不能写成整数方程”，
而是“能不能写成在各个素数处行为良好的整数方程”。
这就是好约化与半稳定约化要解决的问题。

\textbf{核心思想是什么？}
核心思想分两层。
对单条椭圆曲线而言，
我们希望找到一个尽量好的整数模型，
使得模 $p$ 后的纤维要么仍然光滑，
要么只出现最温和的节点奇点；
前者称为好约化，后者称为半稳定约化。
对模曲线而言，
我们则希望整个模问题在整数上也有意义。
这迫使我们把“椭圆曲线”稍微放宽成“广义椭圆曲线”，
从而把尖点也纳入统一的几何框架里。

\textbf{引入之后能得到什么？}
一旦整模型建立起来，
我们就能把 $X_0(N)$ 与 Hecke 对应一起约化到 $\F_p$ 上，
在特殊纤维里真正看见 Frobenius 的几何作用。
对 $p\nmid N$，深刻的 Deligne--Rapoport 与 Igusa 理论说明：
$X_0(N)$ 在 $p$ 处具有好约化；
对 $p\mid N$，特殊纤维虽然会退化，但仍保持半稳定。
这正是下一章把 Hecke 特征值改写成点数公式的根基。

\begin{figure}[ht]
\centering
\begin{tikzpicture}[font=\small, >=Stealth]
  \begin{scope}[xshift=-3.2cm]
    \draw[thick, blue!70!black] plot[smooth cycle] coordinates {(-1.4,0.0) (-0.5,0.9) (0.8,0.7) (1.2,-0.1) (0.2,-0.9) (-1.0,-0.8)};
    \node at (0,-1.5) {好约化：特殊纤维光滑};
  \end{scope}

  \begin{scope}[xshift=3.2cm]
    \draw[thick, orange!80!black] (-1.3,-0.9) .. controls (-0.5,0.4) and (0.5,0.4) .. (1.3,-0.9);
    \draw[thick, orange!80!black] (-1.3,0.9) .. controls (-0.5,-0.4) and (0.5,-0.4) .. (1.3,0.9);
    \filldraw[orange!80!black] (0,0) circle (1.4pt);
    \node at (0,-1.5) {半稳定约化：出现节点};
  \end{scope}

  \draw[->, thick] (-0.8,1.5) -- (0.8,1.5);
  \node at (0,1.85) {模 $p$ 约化};
\end{tikzpicture}
\caption{整模型把复解析对象带到模 $p$ 的世界：当特殊纤维仍然光滑时称为好约化；当只出现普通双点时称为半稳定约化。}
\label{fig:ch07-modular-interpretation}
\end{figure}


先从单条椭圆曲线的判别准则开始。

\begin{definition}[好约化与半稳定约化]
\label{def:ch07-good-semistable-reduction}
设 $E/\Q$ 是椭圆曲线，$p$ 为素数。
若存在 $\Z_p$ 上的模型，其特殊纤维是一条光滑三次曲线，
则称 $E$ 在 $p$ 处有\textbf{好约化}\index{好约化!good reduction}。
若特殊纤维虽不光滑，但只含普通双点（节点）而无尖点，
则称 $E$ 在 $p$ 处有\textbf{半稳定约化}\index{半稳定约化!semistable reduction}。
\end{definition}

\begin{proposition}[判别式判定好约化]
\label{prop:ch07-discriminant-good-reduction}
设
\[
  E: y^2+a_1xy+a_3y=x^3+a_2x^2+a_4x+a_6
\]
是 $\Z_p$ 上的 Weierstrass 方程，判别式为 $\Delta(E)$。
若 $\Delta(E)$ 在 $\Z_p$ 中是单位，
则模 $p$ 的特殊纤维光滑，
因而 $E$ 在 $p$ 处有好约化。
反过来，若模 $p$ 的特殊纤维有奇点，
则 $p\mid \Delta(E)$。
\end{proposition}

\begin{proof}
Weierstrass 方程给出一条平面三次曲线
\[
  F(x,y)=y^2+a_1xy+a_3y-x^3-a_2x^2-a_4x-a_6=0.
\]
模 $p$ 后出现奇点，当且仅当存在点 $(x,y)$ 同时满足
\[
  F(x,y)=0,
  \qquad \frac{\partial F}{\partial x}(x,y)=0,
  \qquad \frac{\partial F}{\partial y}(x,y)=0
\]
于剩余域中成立。
而对 Weierstrass 三次来说，
“方程与两偏导有公共零点”恰好等价于其判别式为零；
这正是判别式的定义性性质：
它刻画多项式出现重根、从而曲线出现奇点的条件。

因此，若 $\Delta(E)$ 在 $\Z_p$ 中为单位，
则其模 $p$ 像非零，特殊纤维不可能有奇点，所以光滑。
反过来，若特殊纤维有奇点，
则模 $p$ 判别式必为零，也就是 $p\mid \Delta(E)$。
证毕。
\end{proof}

\textbf{注。}
对椭圆曲线来说，
真正精确的说法是：
取最小 Weierstrass 模型后，$p\nmid \Delta_{\min}$ 当且仅当好约化；
若 $p\mid \Delta_{\min}$ 但特殊纤维只是节点而不是尖点，
则是乘法型坏约化，也就是半稳定约化。
本章只用到这个判别思想，不展开 Tate 算法。

模曲线的整模型比单条椭圆曲线更微妙，
因为尖点对应的不是光滑椭圆曲线，
而是退化极限。
Deligne--Rapoport 的做法是：
在整数环上不只参数化椭圆曲线，
还允许参数化\textbf{广义椭圆曲线}\index{广义椭圆曲线!generalized elliptic curve}，
也就是在坏纤维里允许出现 N\'eron $n$-gon 那样的节点型退化。
于是尖点不再是“额外补上去的解析边界”，
而是整模型中的真实几何点。

\begin{definition}[整模型与 N\'eron 模型的动机]
\label{def:ch07-integral-neron}
对模曲线 $X_0(N)$，
一个\textbf{整模型}\index{整模型!integral model} 指的是定义在某个整数环（通常是 $\Z$ 或 $\Z[1/N]$）上的曲线
$\mathcal X_0(N)$，
其 generic fiber 为前面得到的 $X_0(N)/\Q$。
对椭圆曲线或 Abel 簇，
若进一步要求群结构在所有光滑测试对象上满足映射延拓的泛性质，
就得到\textbf{N\'eron 模型}\index{N\'eron模型!N\'eron model}。
\end{definition}

本章不证明 Deligne--Rapoport 与 N\'eron 模型的一般存在定理，
只记住它们解决的障碍：
它们让“群结构”“同源”“Hecke 对应”这些本来只在 generic fiber 上可见的对象，
能够一路延伸到整数环与特殊纤维中去。

\begin{example}[$X_0(11)$ 在小素数处的约化]
\label{ex:ch07-x011-reduction}
对 $X_0(11)$，取经典最小模型
\[
  E: y^2+y=x^3-x^2-10x-20.
\]
其判别式为
\[
  \Delta(E)=-11^5.
\]
因此对 $p=2,3,5$，判别式在 $\Z_p$ 中都是单位，
故 $E$ 在这些素数处有好约化。
相反，在 $p=11$ 处判别式被 $11$ 整除，
所以此处必是坏约化；
事实上这里出现的是一个节点，因而属于半稳定约化。
\end{example}

\begin{proof}
对给定方程，按 Weierstrass 不变量公式计算得
\[
  \Delta(E)=-11^5.
\]
于是当 $p=2,3,5$ 时，$p\nmid\Delta(E)$，
由 \cref{prop:ch07-discriminant-good-reduction} 立得好约化。

当 $p=11$ 时，把方程模 $11$ 化简为
\[
  y^2+y=x^3-x^2+x+2.
\]
记
\[
  F(x,y)=y^2+y-x^3+x^2-x-2.
\]
则
\[
  \frac{\partial F}{\partial x}=-3x^2+2x-1,
  \qquad
  \frac{\partial F}{\partial y}=2y+1.
\]
在 $\F_{11}$ 中取 $x=5,y=5$，有
\[
  F(5,5)=0,
  \qquad
  \frac{\partial F}{\partial x}(5,5)=0,
  \qquad
  \frac{\partial F}{\partial y}(5,5)=0.
\]
所以特殊纤维确有奇点。
另一方面，该最小模型对应的导子正是 $11$，
坏约化是乘法型而非加法型，
因此这里出现的是节点而不是尖点，
也就是半稳定约化。
\end{proof}

\textbf{注。}
上例已经把第\ref{ch:jacobians}章的“模椭圆曲线”推进到了整数层面：
同一条曲线在复数域上给出模形式的 Jacobian 因子，
在 $\Q$ 上给出椭圆曲线，
在模 $11$ 后则出现一个节点纤维。
这三幅图像不是彼此分离的，
而是同一个整模型在不同底域上的投影。

对一般模曲线，最重要的深结果可以概括成下面一句话：
Deligne--Rapoport 构造了 $X_0(N)$ 的整模型，
它参数化带 $\Gamma_0(N)$-结构的广义椭圆曲线；
Igusa 进一步证明，当 $p\nmid N$ 时，
该整模型在 $p$ 处的特殊纤维是光滑的。
因此 $X_0(N)$ 在所有不整除级别的素数处都有好约化。
而在 $p\mid N$ 时，特殊纤维则分裂成若干光滑分支并以节点相交，
这就是半稳定约化的几何来源。

这正是下一章最关键的起点：
当 $p\nmid N$ 时，
既然 $X_0(N)$ 在 $\F_p$ 上仍然是一条光滑曲线，
那么 Frobenius 自同态就成为一个真正的代数几何对象；
而上一节的 Hecke 对应也同样可以约化到 $\F_p$ 上。
Eichler--Shimura 关系要做的，
就是比较这两种看似来自不同世界的自同态。

\begin{definition}[广义椭圆曲线与 N\'eron $n$-gon]
\label{def:generalized-elliptic-curve}
所谓\textbf{广义椭圆曲线}，是指一条 proper、flat 的 genus $1$ 曲线，
其光滑部分带有群结构，且几何纤维要么是光滑椭圆曲线，
要么是一个 \textbf{N\'eron $n$-gon}：
也就是 $n$ 个 $\mathbb P^1$ 首尾相接围成的多边形型节点曲线。
这种对象恰好把尖点处出现的退化纤维纳入模问题内部。
\end{definition}

\begin{theorem}[Katz--Mazur，Deligne--Rapoport]
\label{thm:katz-mazur-x0n}
对每个正整数 $N$，$X_0(N)$ 在 $\Z[1/N]$ 上的紧化整模型可以通过广义椭圆曲线及其 $\Gamma_0(N)$-级别结构来描述。
更准确地说，带一个阶 $N$ 循环子群的广义椭圆曲线所组成的模问题，
由一个 proper、flat 的代数曲线表示，
其 generic fiber 正是 $X_0(N)/\Q$。
\end{theorem}

\begin{proof}[说明]
Katz--Mazur 的核心思想是：
不要把尖点视作“解析上补进来的边界”，
而是把它们解释成广义椭圆曲线的真实退化纤维。
这样一来，级别结构、同源与 Hecke 对应都能在整数底上继续定义。
Deligne--Rapoport 最早给出了这一整模型的构造；
Katz--Mazur 则把它整理成系统的模理论语言。
\end{proof}

\begin{proposition}[Tate 曲线是尖点处的局部模型]
\label{prop:tate-curve-generalized-elliptic}
第4章出现的 Tate 曲线 $E_q$ 在形式圆盘 $\operatorname{Spec}\Z[[q]]$ 上可延拓为一条广义椭圆曲线。
其 generic fiber 是通常的椭圆曲线，
而特殊纤维 $q=0$ 是一个 N\'eron $1$-gon，
也就是尖点处最基本的节点纤维。
\end{proposition}

\begin{proof}[说明]
当 $q\neq 0$ 时，Tate 曲线由乘法型统一化 $\mathbb G_m/q^\Z$ 给出，
是光滑椭圆曲线。
把它放到形式邻域里以后，$q=0$ 处的极限不再光滑，
而是恰好退化成一个单环的节点曲线；
这正是广义椭圆曲线定义里允许的最简单坏纤维。
因此尖点附近的 $q$-参数与整模型里的退化几何完全一致。
\end{proof}

\textbf{Deligne--Rapoport 定理的最简表述。}
模堆栈 $\overline{\mathcal M}_{1,1}$ 的紧化可以由广义椭圆曲线来实现：
内部点参数化真正的椭圆曲线，
边界点参数化节点型退化纤维。
对本书而言，最重要的结论并不是堆栈技术本身，
而是：一旦采用广义椭圆曲线的语言，
尖点、整模型、Hecke 对应与模 $p$ 约化就全都变成同一个模问题的不同纤维。
